\documentclass[11pt]{article}
\usepackage[a4paper,margin=27mm]{geometry}
\usepackage{amsmath,amssymb,mathtools}
\usepackage{fontspec,unicode-math}
\setmainfont{Latin Modern Roman}
\setmathfont{Latin Modern Math}
\usepackage{xurl,hyperref}
\hypersetup{colorlinks=true,urlcolor=blue,linkcolor=blue}
\newcommand{\sref}[2]{\hyperlink{source:#1:#2}{[S#2]}}
\newcommand{\cataloguescope}[1]{\paragraph{Mathematical question.} #1}
\begin{document}
\section{Separatrix Splitting under Quasiperiodic Forcing}\label{statement}

\subsection{Definitions and mathematical statement}

\paragraph{Shared notation.}$\mathbb N=\{1,2,\ldots\}$, and $\mathbb Z,\mathbb Q,\mathbb R,\mathbb C$ denote the integers, rationals, reals and complex numbers. Any other conventions are specified locally.


Fix $J,g>0$ and even zero-average trigonometric polynomials $f$ on $(\mathbb R/2\pi\mathbb Z)^2$ and $m$ on $\mathbb R/2\pi\mathbb Z$. For actions $A\in\mathbb R^2,I\in\mathbb R$ and angles $\alpha\in(\mathbb R/2\pi\mathbb Z)^2,\varphi\in\mathbb R/2\pi\mathbb Z$, set \[H=\omega(\eta)\cdot A+\frac{I^2}{2J}+Jg^2(\cos\varphi-1)+\varepsilon f(\alpha)(\cos\varphi-1)+\varepsilon^{\prime}m(2\alpha_2)(\cos\varphi-1),\] where $\omega(\eta)=(\eta^{1/2}\overline\omega_1,\eta^{-1/2}\overline\omega_2)$. Let $\eta_n\downarrow0$ satisfy $|\omega(\eta_n)\cdot\nu|>\eta_n^a|\nu|^{-\tau}$ for every $0\ne\nu\in\mathbb Z^2$, with fixed $a,\tau>0$. The invariant torus has $A=0,I=0,\varphi=0$. On the section $H=0,\varphi=\pi$, write its unstable and stable manifolds as graphs $A^u(\alpha),A^s(\alpha)$ when these exist and put $\Delta(\varepsilon,\varepsilon^{\prime},\eta)=\det[\partial_\alpha(A^u-A^s)]_{\alpha=0}$. The evenness gives the indicated symmetric homoclinic point.
\paragraph{Statement recorded in the supplied source.}
With $\varepsilon=\eta^c$ for a fixed sufficiently large $c>0$, determine the asymptotic splitting $\Delta(\eta^c,\varepsilon^{\prime},\eta)$ as $\eta=\eta_n\to0$ when $\varepsilon^{\prime}=O(1)$, or at least obtain a positive lower bound for its absolute value in the nondegenerate cases. This includes understanding the transition in the splitting asymptotics when $\varepsilon^{\prime}=\eta^T$ crosses a critical exponent $T$, despite convergence of the perturbation series throughout this larger parameter region. \sref{statement}{2}
\subsection{Short English statement}
Find the exponentially small separatrix splitting when one quasiperiodic forcing component is small and the fast component remains of bounded size, including the change in asymptotics inside the perturbation series convergence region.
\subsection{Sources}
\begin{description}
\item[\hypertarget{source:statement:1}{S1}] ULAM.AI and the UnsolvedMath Contributors, UnsolvedMath: A Curated Collection of Open Mathematics Problems, catalogue entry. CC BY 4.0. \url{https://huggingface.co/datasets/ulamai/UnsolvedMath}.
\item[\hypertarget{source:statement:2}{S2}] Giovanni Gallavotti, Separatrix Splitting for a Quasi Periodically Forced Pendulum (6 November 1998), Hamiltonian, definition of Delta, Diophantine sequence, and final Question. \url{https://web.math.princeton.edu/~aizenman/OpenProblems_MathPhys/9805.separatrix.tex}.
\end{description}
\label{end:statement}
\end{document}
